\chapter{Détermination du barycentre}

Soit un graphe $\mathcal{G}$, comme expliqué en \ref{partie:graph},  composé d’un ensemble de nœuds, noté $\mathcal{E}$, et d’un ensemble d’arêtes, noté $\mathcal{V}$.
\begin{itemize}
    \item $\mathcal{V}=\{1, 2, \ldots, n\}$, l'ensemble des noeuds.
    \item $\mathcal{E}\subseteq \{(i, j) : i, j \in \mathcal{V},\ j \ne i\}$. Une arête est composée de noeuds $i$ et $j$
\end{itemize}

\section{Définition }

Soit un ensemble de $n$ robots dont les positions respectives à l’instant $t$ sont notées $p_i(t) \in \mathbb{R}^d$, avec $d = 2$ pour un plan ou $d = 3$ pour un espace tridimensionnel.
Le \textbf{barycentre} de l’essaim est défini comme la moyenne arithmétique des positions des robots :

\begin{equation}
p_{\text{bary}}(t) = \frac{1}{n} \sum_{i=1}^{n} p_i(t)
\end{equation}

Il s’agit du centre géométrique de l’ensemble des robots à un instant donné.

\section{Calcul du barycentre au sein d'un essaim distribué }

Dans le cas d’un essaim de robots distribué, chaque robot n'a accès qu'à une quantité limité d'information : il peut uniquement communiquer avec ses voisins immédiats, définis par les arêtes du graphe $\mathcal{G}$.  
Le calcul du barycentre global de l’essaim nécessite donc des échanges d’information entre les robots. Plusieurs méthodes existent pour estimer ce barycentre de manière distribuée.

\subsection{Méthode flooding}

La méthode dite \textit{flooding} consiste à propager l’information brute (les positions individuelles des robots) à travers le graphe de communication $\mathcal{G}$, de manière itérative.  
L'objectif est que chaque robot puisse collecter la position de tous les robots de l’essaim, même ceux qui ne sont pas directement connectés à lui.
Le déroulement de la méthode est le suivant :

\begin{itemize}
    \item À l’instant initial, chaque robot $i$ connaît uniquement sa propre position $p_i(t)$.
    \item À la première étape, chaque robot envoie sa position $p_i(t)$ à ses voisins directs $\mathcal{N}_i$.
    \item À l’étape suivante, chaque robot transmet à ses voisins non seulement sa propre position $p_i(t)$, mais aussi toutes les positions qu’il a reçues des autres robots au cours des étapes précédentes.
    \item Ce processus se poursuit jusqu’à ce que toutes les positions $\{ p_j(t) \}_{j=1}^{n}$ aient circulé sur l’ensemble du graphe, et que chaque robot dispose de l’information complète.
\end{itemize}

À l’issue de ce processus, chaque robot peut calculer localement et exactement le barycentre de l’essaim :

\begin{equation}
p_{\text{bary}}(t) = \frac{1}{n} \sum_{j=1}^{n} p_j(t)
\end{equation}

Cette méthode assure un résultat exact et ne dépend pas du nombre d’itérations.  
Cependant, elle est coûteuse en communication, car le volume de données augmente rapidement avec le nombre de robots — chaque robot doit relayer un vecteur contenant les positions de tous les robots qu’il connaît déjà.
Le nombre total de messages échangés dans le réseau est de l’ordre de :

\begin{equation}
\mathcal{O}(n \cdot \mathcal{E})
\end{equation}

où $n$ est le nombre de robots et $\mathcal{E}$ est le nombre d’arêtes (liens de communication) du graphe $\mathcal{G}$.

Chaque position de robot doit en effet être relayée à travers le graphe, en empruntant potentiellement plusieurs arêtes. Le coût en communication augmente donc linéairement avec le nombre de robots $n$, et proportionnellement à la taille du graphe.

\subsection{Méthode de consensus}

Une alternative plus légère en communication consiste à utiliser une \textit{méthode de consensus}.  
Dans cette approche, chaque robot initialise une estimation locale de la position du barycentre :

\begin{equation}
p_i^{(0)} = p_i(t)
\end{equation}

À chaque itération $k$, le robot $i$ échange son estimation $p_i^{(k)}$ avec ses voisins, puis met à jour son estimation selon :

\begin{equation}
p_i^{(k+1)} = w_{ii} p_i^{(k)} + \sum_{j \in \mathcal{N}_i} w_{ij} p_j^{(k)}
\end{equation}

où $\mathcal{N}_i$ est l’ensemble des voisins de $i$, et $w_{ij}$ sont des coefficients de pondération tels que :

\begin{equation}
\sum_{j \in \mathcal{N}_i \cup \{i\}} w_{ij} = 1, \quad w_{ij} \geq 0
\end{equation}

Par exemple, on peut utiliser des poids uniformes :

\begin{equation}
w_{ij} = \frac{1}{d_i + 1}
\end{equation}

avec $d_i$ le degré (nombre de voisins) du nœud $i$.

---

Au fil des itérations, les estimations $p_i^{(k)}$ de chaque robot convergent vers le barycentre réel $p_{\text{bary}}(t)$, sous réserve que le graphe soit connexe et que les coefficients $w_{ij}$ soient bien choisis.

Cette méthode présente plusieurs avantages :
\begin{itemize}
    \item Communication uniquement locale (entre voisins).
    \item Réduction du coût de communication par rapport au flooding.
    \item Estimation progressive : le nombre d’itérations $k$ permet de gérer le compromis précision/communication.
\end{itemize}

En contrepartie, la convergence vers le barycentre exact n’est jamais instantanée et dépend de la structure du graphe et du nombre d’itérations réalisées.